% Л4 -- C#: constructs, patterns, delegates and errors. Slides.
%
% Companion to the spoken script in ../L04-csharp-constructs.md, which is
% the source of truth: the slides carry code, diagrams and short claims, and
% the delivery lives in the script only: by the author's decision this deck
% carries no speaker notes. Slide text is in Russian because that is the
% content; everything else in this file is English.
%
% Build:  make          (latexmk -pdf; engine: pdflatex)
%
% Every listing was compiled on .NET 10; outputs in comments are real.
% Frame 48 is real terminal output, with the path anonymised to /home/user.
%
% House rules: ../CLAUDE.md and ../../../CLAUDE.md. Theme: ../../../../theme/.

\documentclass[aspectratio=169,10pt,t]{beamer}

\usetheme{dotnet}
\usetikzlibrary{arrows.meta,positioning,calc}
\usepackage{array}   % ragged-right table columns

\dotnetlecture{Л4}

\title[Л4. Конструкции C\#]{C\#: конструкции, делегаты и ошибки}
\author{Дмитрий Курочкин}
\date{}

% A listing set one step under \scriptsize, for two columns of code side by
% side. Smaller than this does not carry to the back of the room.
\newcommand{\twocol}{\ttfamily\fontsize{7.0}{8.2}\selectfont}

\begin{document}

% ------------------------------------------------------------------ title --
\begin{frame}[plain]
  \titlepagednc{Лекция 4}{C\#: конструкции, делегаты и ошибки}{}{Дмитрий
    Курочкин}
\end{frame}

% ======================================================= 1. CONSTRUCTS =====
\section{1. Управляющие конструкции}

\begin{frame}[plain]
  \sectionopener{1}{Управляющие конструкции}{}
\end{frame}

\begin{frame}{Лестница ветвлений}
  \slidetop
  \claim{От простого к сложному: каждая ступень --- та же задача, записанная
    \hi{короче}.}

  \begin{center}
    \begin{tikzpicture}[font=\small,
      st/.style={draw=hairline,line width=0.8pt,fill=codebg,
                 minimum height=0.95cm,minimum width=2.2cm,align=center,
                 inner sep=0.3em},
      ar/.style={-{Straight Barb[length=1.8mm]},muted,line width=0.6pt}]
      \node[st,draw=alt]    (a) at (0,0)    {\texttt{if}};
      \node[st,draw=accent] (b) at (2.75,0) {\texttt{? :}};
      \node[st,draw=alt]    (c) at (5.5,0)  {\texttt{switch}};
      \node[st,draw=accent] (d) at (8.25,0) {\texttt{switch}\\[-0.1em]
        {\scriptsize выражение}};
      \node[st,dashed,fill=paper] (e) at (11.0,0) {образцы\\[-0.1em]
        {\scriptsize раздел 2}};
      \draw[ar] (a) -- (b); \draw[ar] (b) -- (c);
      \draw[ar] (c) -- (d); \draw[ar] (d) -- (e);
      \foreach \n in {a,c} \node[above=0.35em,font=\scriptsize\bfseries,
        text=alt] at (\n.north) {инструкция};
      \foreach \n in {b,d} \node[above=0.35em,font=\scriptsize\bfseries,
        text=accent] at (\n.north) {выражение};
      \foreach \n/\t in {a/два исхода,b/два исхода,c/много исходов,
                         d/много исходов} \node[below=0.35em,
        font=\scriptsize,text=muted] at (\n.south) {\t};
    \end{tikzpicture}
  \end{center}

  \slidegap
  \aside{\lo{Инструкция} выполняет действие, результат собирают через
    переменную. \hi{Выражение} возвращает значение: его присваивают,
    возвращают, передают.}
\end{frame}

\begin{frame}[fragile]{\texttt{if}: ветвление по условию}
  \slidetop
  \begin{minipage}[t]{0.46\textwidth}
\begin{lstlisting}[style=csharp,basicstyle=\ttfamily\small]
decimal shipping;
if (order.Total >= 5000m)
{
    shipping = 0m;
}
else
{
    shipping = 300m;
}
\end{lstlisting}
  \end{minipage}\hfill
  \begin{minipage}[t]{0.48\textwidth}
    \vspace{0.2em}
    \begin{cbul}{accent}
      \item условие --- всегда \hi{\texttt{bool}}
      \item ветка выполнится, если условие истинно
      \item \texttt{else} --- на случай лжи, необязателен
    \end{cbul}
    \vspace{1.2em}
    \aside{Обе ветки делают одно и то же: присваивают \texttt{shipping},
      объявленную \lo{до} \texttt{if}. Запомните этот листинг.}
  \end{minipage}
\end{frame}

\begin{frame}[fragile]{Лестница: плоская, а не вложенная}
  \slidetop
  \begin{minipage}[t]{0.47\textwidth}
    \colhead{alt}{вложенная}\par\vspace{0.5em}
\begin{lstlisting}[style=csharp,basicstyle=\ttfamily\scriptsize]
if (score >= 90) {
    grade = "A";
} else {
    if (score >= 75) {
        grade = "B";
    } else {
        if (score >= 60) {
            grade = "C";
        } else {
            grade = "F";
        }
    }
}
\end{lstlisting}
  \end{minipage}\hfill
  \begin{minipage}[t]{0.47\textwidth}
    \colhead{accent}{плоская}\par\vspace{0.5em}
\begin{lstlisting}[style=csharp,basicstyle=\ttfamily\scriptsize]
if (score >= 90) {
    grade = "A";
} else if (score >= 75) {
    grade = "B";
} else if (score >= 60) {
    grade = "C";
} else {
    grade = "F";
}
\end{lstlisting}
    \vspace{0.8em}
    \aside{Сверху вниз, срабатывает \hi{первое} истинное условие;
      \texttt{else} ловит остальное. Частное --- выше общего.}
  \end{minipage}
\end{frame}

\begin{frame}[fragile]{Ранний выход}
  \slidetop
  \begin{minipage}[t]{0.45\textwidth}
    \colhead{alt}{всё внутри if}\par\vspace{0.5em}
\begin{lstlisting}[style=csharp,basicstyle=\ttfamily\fontsize{6.8}{7.9}\selectfont]
string Pay(Order? order, User user) {
    if (order is not null) {
        if (user.CanPay) {
            if (order.IsNew) {
                order.MarkPaid();
                return "paid";
            } else {
                return "already paid";
            }
        } else {
            return "forbidden";
        }
    } else {
        return "no order";
    }
}
\end{lstlisting}
  \end{minipage}\hfill
  \begin{minipage}[t]{0.51\textwidth}
    \colhead{accent}{проверки сверху}\par\vspace{0.5em}
\begin{lstlisting}[style=csharp,basicstyle=\ttfamily\fontsize{6.8}{7.9}\selectfont]
string Pay(Order? order, User user) {
    if (order is null) return "no order";
    if (!user.CanPay) return "forbidden";
    if (!order.IsNew) return "already paid";

    order.MarkPaid();
    return "paid";
}
\end{lstlisting}
    \vspace{0.8em}
    \aside{Список условий, потом суть --- на \hi{первом} уровне отступа,
      без единого \texttt{else}. Слева, чтобы понять строку на четвёртом
      уровне, надо помнить \lo{все условия} над ней.}
  \end{minipage}
\end{frame}

\begin{frame}[fragile]{Тернарный оператор}
  \slidetop
  \claim{\texttt{if}, который \hi{возвращает значение}.}

  \begin{minipage}[t]{0.34\textwidth}
    \colhead{alt}{инструкция}\par\vspace{0.5em}
\begin{lstlisting}[style=csharp,basicstyle=\ttfamily\scriptsize]
decimal shipping;
if (order.Total >= 5000m)
{
    shipping = 0m;
}
else
{
    shipping = 300m;
}
\end{lstlisting}
  \end{minipage}\hfill
  \begin{minipage}[t]{0.61\textwidth}
    \colhead{accent}{выражение}\par\vspace{0.5em}
\begin{lstlisting}[style=csharp,basicstyle=\ttfamily\scriptsize]
var shipping = order.Total >= 5000m ? 0m : 300m;
\end{lstlisting}
    \vspace{1.0em}
    \aside{Условие, \texttt{?}, значение для истины, \texttt{:}, значение
      для лжи. Результат можно присвоить, вернуть, передать в метод,
      вставить в интерполяцию.}\par
    \vspace{0.8em}
    \aside{Хорош, пока \hi{короткий}. Тернарный внутри тернарного ---
      \lo{сигнал} брать \texttt{switch}.}
  \end{minipage}
\end{frame}

\begin{frame}[fragile]{\texttt{switch}: было / стало}
  \slidetop
  \begin{minipage}[t]{0.485\textwidth}
    \colhead{alt}{инструкция}\par\vspace{0.5em}
\begin{lstlisting}[style=csharp,basicstyle=\ttfamily\fontsize{6.4}{7.4}\selectfont]
string label;
switch (status)
{
    case OrderStatus.New:
        label = "new";
        break;
    case OrderStatus.Paid:
        label = "paid";
        break;
    case OrderStatus.Shipped:
        label = "shipped";
        break;
    default:
        throw new ArgumentOutOfRangeException(
            nameof(status));
}
\end{lstlisting}
  \end{minipage}\hfill
  \begin{minipage}[t]{0.485\textwidth}
    \colhead{accent}{выражение}\par\vspace{0.5em}
\begin{lstlisting}[style=csharp,basicstyle=\ttfamily\fontsize{6.4}{7.4}\selectfont]
var label = status switch
{
    OrderStatus.New => "new",
    OrderStatus.Paid => "paid",
    OrderStatus.Shipped => "shipped",
    _ => throw new ArgumentOutOfRangeException(
        nameof(status)),
};
\end{lstlisting}
    \vspace{0.8em}
    \aside{\texttt{switch} после значения, стрелки вместо \texttt{:} и
      \texttt{break}, \texttt{\_} вместо \texttt{default}. Выражение
      \hi{обязано} вернуть значение.}
  \end{minipage}
\end{frame}

\begin{frame}[fragile]{\texttt{when}: условие к ветке}
  \slidetop
  \begin{center}
    \begin{tikzpicture}[font=\ttfamily\small,
      part/.style={draw=hairline,line width=0.8pt,fill=codebg,
                   minimum height=0.75cm,inner sep=0.45em},
      lab/.style={font=\normalfont\scriptsize,text=muted,align=center,
                  text width=3.4cm,below=0.3em}]
      \node[part] (p) at (0,0) {var s};
      \node[part,draw=accent,line width=1.0pt,anchor=west]
        (w) at ([xshift=1.1cm]p.east) {\hi{when} s >= 90};
      \node[part,anchor=west] (r) at ([xshift=1.1cm]w.east) {=> "A"};
      \node[lab] at (p.south) {образец: любое число,\\ назовём его \texttt{s}};
      \node[lab] at (w.south) {условие: ветка подходит,\\ только если оно
        \hi{истинно}};
      \node[lab] at (r.south) {результат\\ ветки};
    \end{tikzpicture}
  \end{center}

  \vspace{0.4em}
  \begin{minipage}[t]{0.47\textwidth}
    \colhead{alt}{лестница if}\par\vspace{0.5em}
\begin{lstlisting}[style=csharp,basicstyle=\twocol]
string grade;

if (score == 100) grade = "A+";
else if (score >= 90) grade = "A";
else if (score >= 75) grade = "B";
else if (score >= 60) grade = "C";
else grade = "F";

\end{lstlisting}
  \end{minipage}\hfill
  \begin{minipage}[t]{0.47\textwidth}
    \colhead{accent}{switch с when}\par\vspace{0.5em}
\begin{lstlisting}[style=csharp,basicstyle=\twocol]
var grade = score switch
{
    100 => "A+",
    var s when s >= 90 => "A",
    var s when s >= 75 => "B",
    var s when s >= 60 => "C",
    _ => "F",
};
\end{lstlisting}
  \end{minipage}
\end{frame}

\begin{frame}{Инструкция и выражение}
  \slidetop
  \claim{Switch-выражение --- это \hi{тернарный оператор} в семье
    \texttt{switch}.}

  {\renewcommand{\arraystretch}{1.6}
  \begin{tabular}{@{}p{3.6cm}p{4.2cm}p{4.6cm}@{}}
    & \colhead{alt}{инструкция} & \colhead{accent}{выражение} \\
    \textbf{два исхода}    & \texttt{if} / \texttt{else}
                           & \texttt{\hi{cond ? a : b}} \\
    \textbf{много исходов} & \texttt{switch} / \texttt{case}
                           & \texttt{\hi{x switch \{ ... \}}} \\
  \end{tabular}}

  \slidegap
  \hrulethin{0.42}
  \rulegap
  \aside{Выбор между столбцами один и тот же: нужно \lo{действие} или нужно
    \hi{значение}. Захотелось в ветке записать в лог и вернуть --- это
    инструкция, и честнее взять её.}
\end{frame}

\begin{frame}[plain]
  \vote{Что такое \texttt{var}?}
\end{frame}

\begin{frame}[fragile]{\texttt{var}: вывод типа, а не динамика}
  \slidetop
  \claim{Тип у переменной есть, он один и известен \hi{при компиляции}.}

  \begin{minipage}[t]{0.47\textwidth}
    \colhead{accent}{тип виден справа}\par\vspace{0.5em}
\begin{lstlisting}[style=csharp,basicstyle=\ttfamily\scriptsize]
var customers = new List<Customer>();
var point = new Point(1, 2);
var total = (decimal)row["total"];
\end{lstlisting}
  \end{minipage}\hfill
  \begin{minipage}[t]{0.47\textwidth}
    \colhead{alt}{тип не виден}\par\vspace{0.5em}
\begin{lstlisting}[style=csharp,basicstyle=\ttfamily\scriptsize]
var result = Process(order);   // ?
Invoice result = Process(order);
\end{lstlisting}
  \end{minipage}

  \slidegap
  \hrulethin{0.42}
  \rulegap
  {\large Тип должен быть виден в строке \hi{хотя бы один раз}.}
\end{frame}

\begin{frame}[fragile]{\texttt{new} без имени типа}
  \slidetop
  \claim{Тип известен из контекста --- пишем его \hi{один раз}.}

\begin{lstlisting}[style=csharp,basicstyle=\ttfamily\small]
Dictionary<string, List<int>> scores = new();
private readonly List<Order> _orders = new();
Point origin = new(0, 0);
Draw(new(10, 20));                    // parameter of type Point
\end{lstlisting}

  \slidegap
  \aside{Работает при присваивании в переменную объявленного типа, в поле,
    в параметр и при возврате из метода. То же правило, что у
    \texttt{var}, только \hi{в другую сторону}: тип слева, короткий
    \texttt{new} справа.}
\end{frame}

\begin{frame}[fragile]{Выражения коллекций}
  \slidetop
  \begin{minipage}[t]{0.52\textwidth}
    \colhead{alt}{три синтаксиса}\par\vspace{0.5em}
\begin{lstlisting}[style=csharp,basicstyle=\twocol]
int[] a = new int[] { 1, 2, 3 };
List<string> b = new List<string> { "x", "y" };
Span<int> c = stackalloc int[] { 1, 2, 3 };

var all = new List<int>(a);
all.Add(4);
var none = new List<int>();
\end{lstlisting}
  \end{minipage}\hfill
  \begin{minipage}[t]{0.44\textwidth}
    \colhead{accent}{один}\par\vspace{0.5em}
\begin{lstlisting}[style=csharp,basicstyle=\twocol]
int[] a = [1, 2, 3];
List<string> b = ["x", "y"];
Span<int> c = [1, 2, 3];

List<int> all = [.. a, 4];

List<int> none = [];
\end{lstlisting}
  \end{minipage}

  \slidegap
  \aside{\texttt{\hi{..}} --- оператор расширения: вставить все элементы
    другой коллекции. Представление компилятор выбирает под целевой тип:
    массив, список, Span на стеке. Пустые скобки --- \hi{дешёвая} пустая
    коллекция.}
\end{frame}

\begin{frame}[fragile]{Индексы и диапазоны}
  \slidetop
  \begin{minipage}[t]{0.55\textwidth}
    \colhead{alt}{арифметика}\par\vspace{0.5em}
\begin{lstlisting}[style=csharp,basicstyle=\twocol]
var last = items[items.Length - 1];
var beforeLast = items[items.Length - 2];
var middle = items.Skip(1)
    .Take(items.Length - 2).ToArray();
\end{lstlisting}
  \end{minipage}\hfill
  \begin{minipage}[t]{0.41\textwidth}
    \colhead{accent}{намерение}\par\vspace{0.5em}
\begin{lstlisting}[style=csharp,basicstyle=\twocol]
var last = items[^1];
var beforeLast = items[^2];
var middle = items[1..^1];
\end{lstlisting}
  \end{minipage}

  \slidegap
  \aside{<<Последний элемент>> и <<длина минус один>> --- разные
    высказывания. Первое говорит, \hi{что} вы хотите, второе --- \lo{как}
    вычисляете, и в вычислении можно ошибиться на единицу.}
\end{frame}

\begin{frame}{Счёт с начала и с конца}
  \slidetop
  \begin{center}
    \begin{tikzpicture}[font=\small,
      cell/.style={draw=hairline,line width=0.8pt,fill=codebg,
                   minimum width=1.5cm,minimum height=0.9cm,inner sep=0pt}]
      \foreach \v [count=\i from 0] in {10,20,30,40,50} {
        \node[cell] (c\i) at (\i*1.5,0) {\texttt{\v}};
        \node[accent,font=\small\ttfamily,above=0.3em] at (c\i.north) {\i};
      }
      \foreach \i/\l in {0/\textasciicircum5,1/\textasciicircum4,
                         2/\textasciicircum3,3/\textasciicircum2,
                         4/\textasciicircum1} {
        \node[alt,font=\small\ttfamily,below=0.3em] at (c\i.south) {\l};
      }
      \draw[accent,line width=1.6pt] ([xshift=1pt,yshift=1pt]c1.south west)
        rectangle ([xshift=-1pt,yshift=-1pt]c3.north east);
      \node[accent,font=\small\bfseries,anchor=west] at (8.0,0)
        {\texttt{1..\textasciicircum1}};
      \node[muted,font=\scriptsize,anchor=west] at (8.0,-0.45)
        {= 20, 30, 40};
    \end{tikzpicture}
  \end{center}

  \slidegap
  {\small
  \begin{tabular}{@{}p{2.6cm}p{6.0cm}@{}}
    \texttt{\hi{..\textasciicircum1}} & всё, кроме последнего \\[0.2em]
    \texttt{\hi{1..}}  & всё, начиная со второго \\[0.2em]
    \texttt{\hi{1..\textasciicircum1}} & без первого и последнего \\
  \end{tabular}}
\end{frame}

\begin{frame}[fragile]{Диапазон: копия или окно}
  \slidetop
\begin{lstlisting}[style=csharp,basicstyle=\ttfamily\small]
Range body = 1..^1;                         // a value like any other
int[] copy = array[body];                   // new array
Span<int> window = array.AsSpan()[body];    // the same memory
\end{lstlisting}

  \slidegap
  \versus{на Span --- окно}{%
    \begin{cbul}{accent}
      \item та же память, без копирования
      \item режьте Span в горячем пути
    \end{cbul}}{на массиве --- копия}{%
    \begin{cbul}{alt}
      \item новый массив, копия куска
      \item синтаксис тот же, стоимость другая
    \end{cbul}}

  \slidegap
  \aside{\texttt{Index} и \texttt{Range} --- обычные типы стандартной
    библиотеки: их можно сохранить в переменную и передать в метод.}
\end{frame}

\begin{frame}[plain]
  \keyline{Пишите выражениями, а не инструкциями:\\ выражение обязано
    вернуть значение, и это делает его честнее.}
\end{frame}

% ================================================= 2. PATTERN MATCHING =====
\section{2. Pattern matching}

\begin{frame}[plain]
  \sectionopener{2}{Pattern matching}{Шесть ступеней: константа $\rightarrow$
    тип $\rightarrow$ свойство $\rightarrow$ отношение и логика
    $\rightarrow$ позиция $\rightarrow$ \hi{список}.}
\end{frame}

\begin{frame}[fragile]{Одно выражение вместо цепочки приведений}
  \slidetop
\begin{lstlisting}[style=csharp,basicstyle=\ttfamily\fontsize{7.2}{8.3}\selectfont]
static string Describe(object? input) => input switch
{
    null => "nothing",
    int and < 0 => "negative number",
    int n => $"number {n}",
    string { Length: 0 } => "empty string",
    string s => $"string of {s.Length} chars",
    Point { X: 0, Y: 0 } => "origin",
    Point(var x, var y) when x == y => $"on the diagonal: {x}",
    Point p => $"point {p.X}:{p.Y}",
    int[] and [] => "empty array",
    int[] and [var first, ..] => $"array, first is {first}",
    _ => "something else",
};
\end{lstlisting}
  \vspace{0.3em}
  \aside{Проверка и извлечение \hi{одновременно}: ни одного каста, ни одной
    временной переменной, ни одного места, где можно проверить одно, а
    привести \lo{к другому}.}
\end{frame}

\begin{frame}[fragile]{Pattern matching в \texttt{if}}
  \slidetop
  \begin{minipage}[t]{0.5\textwidth}
\begin{lstlisting}[style=csharp,basicstyle=\ttfamily\small]
if (input is string s)
{
    Console.WriteLine(s.Length);
}

if (input is not string text)
{
    return;
}
Console.WriteLine(text.Length);
\end{lstlisting}
  \end{minipage}\hfill
  \begin{minipage}[t]{0.44\textwidth}
    \vspace{0.2em}
    \begin{cbul}{accent}
      \item внутри первого \texttt{if} \texttt{s} уже строка
      \item \texttt{\hi{is not}} и выход: ниже \texttt{text} --- строка
            \hi{во всём методе}
    \end{cbul}
    \vspace{1.0em}
    \aside{Проверка, приведение и ранний выход из первого раздела --- одной
      строкой. Так начинается половина методов, разбирающих вход.}
  \end{minipage}
\end{frame}

\begin{frame}[fragile]{Tuple: вернуть два значения}
  \slidetop
\begin{lstlisting}[style=csharp,basicstyle=\ttfamily\small]
var (lo, hi) = MinMax([5, 2, 9]);
Console.WriteLine($"{lo}..{hi}");   // 2..9

static (int Min, int Max) MinMax(IEnumerable<int> values)
{
    var min = int.MaxValue;
    var max = int.MinValue;
    foreach (var v in values)
    {
        if (v < min) min = v;
        if (v > max) max = v;
    }
    return (min, max);
}
\end{lstlisting}
\end{frame}

\begin{frame}{Где уместен tuple}
  \slidetop
  \claim{Tuple \hi{локальный}, record \lo{публичный}.}

  \versus{tuple}{%
    \begin{cbul}{accent}
      \item вернуть из метода и тут же разобрать
      \item два-три значения
      \item вместо \texttt{out}-параметров
    \end{cbul}}{record}{%
    \begin{cbul}{alt}
      \item параметр метода
      \item поле, коллекция, наружу из класса
      \item данные, которые живут дольше вызова
    \end{cbul}}

  \slidegap
  \aside{Tuple --- структура, сравнивается по значению. Имена элементов есть
    только для компилятора: во время выполнения их нет.}
\end{frame}

\begin{frame}[fragile]{Деконструкция}
  \slidetop
\begin{lstlisting}[style=csharp,basicstyle=\ttfamily\small]
var (_, max) = MinMax(values);      // only Max is needed
(a, b) = (b, a);                    // swap, no temporary
var (x, y) = new Point(3, 4);       // a record deconstructs too
if ((x, y) == (0, 0)) { }           // element-wise comparison
\end{lstlisting}

  \slidegap
  \aside{Любой тип разбирается, если у него есть метод \texttt{Deconstruct}.
    Поэтому tuple и pattern matching --- \hi{один механизм}: то, что можно разобрать,
    можно и сопоставить.}
\end{frame}

\begin{frame}[plain]
  \keyline{Pattern matching переносит проверку формы данных на компилятор:\\
    вы описываете, что ожидаете, а он проверяет и извлекает.}
\end{frame}

% ========================================================= 3. DELEGATES ====
\section{3. Делегаты, лямбды, замыкания}

\begin{frame}[plain]
  \sectionopener{3}{Делегаты, лямбды, замыкания}{Делегат --- это \hi{тип},
    описывающий сигнатуру метода.}
\end{frame}

\begin{frame}{Готовые формы}
  \slidetop
  \claim{Свои делегаты объявляют \hi{редко}: стандартная библиотека уже
    объявила все формы.}

  {\renewcommand{\arraystretch}{1.6}
  \begin{tabular}{@{}p{4.6cm}p{2.4cm}p{6.2cm}@{}}
    \texttt{\hi{Func<T, TResult>}} & \textbf{вычисли} & принимает 0--16
      параметров, последний тип --- результат \\
    \texttt{\hi{Action<T>}}         & \textbf{сделай}  & то же, ничего не
      возвращает \\
    \texttt{Predicate<T>}           & \textbf{проверь} & один параметр,
      \texttt{bool}; в новом коде --- \texttt{Func<T, bool>} \\
  \end{tabular}}

  \slidegap
  \aside{Экземпляр делегата --- пара: ссылка на метод и ссылка на объект,
    у которого его вызывать. Поэтому \lo{не указатель}: указатель не знает
    про объект.}
\end{frame}

\begin{frame}[fragile]{Формы лямбды}
  \slidetop
\begin{lstlisting}[style=csharp,basicstyle=\ttfamily\small]
Func<int, bool> isEven = x => x % 2 == 0;
Func<int, int, int> add = (a, b) => a + b;
Action<string> print = Console.WriteLine;     // method group
Func<string, int> parse = s =>
{
    var n = int.Parse(s);
    return n * 2;
};
Func<Task> wait = async () => await Task.Delay(100);
var square = (int x) => x * x;                // Func<int, int>
\end{lstlisting}

  \slidegap
  \aside{Параметры, стрелка, тело. Одно выражение --- без скобок и
    \texttt{return}. Типы параметров компилятор выводит \hi{из ожидаемого
    делегата}.}
\end{frame}

\begin{frame}[fragile]{Замыкания: захват переменной, а не значения}
  \slidetop
  \begin{minipage}[t]{0.66\textwidth}
\begin{lstlisting}[style=csharp,basicstyle=\ttfamily\fontsize{7.6}{9.0}\selectfont]
var counter = 0;
Action increment = () => counter++;
increment();
increment();
Console.WriteLine(counter);

var actions = new List<Action>();
for (var i = 0; i < 3; i++)
    actions.Add(() => Console.Write(i));
foreach (var a in actions) a();

actions.Clear();
foreach (var j in new[] { 0, 1, 2 })
    actions.Add(() => Console.Write(j));
foreach (var a in actions) a();
\end{lstlisting}
  \end{minipage}\hfill
  \begin{minipage}[t]{0.30\textwidth}
    \vspace{2.2em}
    {\large\ttfamily\hi{2}}\enspace\aside{та же переменная}\par
    \vspace{3.75em}
    {\large\ttfamily\lo{333}}\enspace\aside{одна \texttt{i} на весь цикл}\par
    \vspace{3.75em}
    {\large\ttfamily\hi{012}}\enspace\aside{новая \texttt{j} на итерацию}
  \end{minipage}
\end{frame}

\begin{frame}{Что живёт дольше, чем вы думали}
  \slidetop
  \claim{Захваченное живёт, пока жива \hi{лямбда}, а не метод.}

  \begin{center}
    \begin{tikzpicture}[font=\small,
      b/.style={draw=hairline,line width=0.8pt,fill=codebg,
                minimum height=0.9cm,align=center,inner sep=0.4em},
      ar/.style={-{Straight Barb[length=1.8mm]},line width=0.7pt}]
      \node[b,minimum width=2.6cm] (m) at (0,0) {метод\\[-0.1em]
        {\scriptsize\color{muted} кадр стека}};
      \node[b,draw=alt,minimum width=3.4cm] (h) at (4.6,0)
        {скрытый класс\\[-0.1em]{\scriptsize\ttfamily int counter;
         int[] big;}};
      \node[b,draw=accent,minimum width=2.6cm] (l) at (9.2,0) {лямбда\\[-0.1em]
        {\scriptsize\color{muted} в кэше, в поле, в событии}};
      \draw[ar,muted] (m) -- node[above,font=\scriptsize]{создаёт} (h);
      \draw[ar,accent] (l) -- node[above,font=\scriptsize]{держит} (h);
      \draw[ar,alt,dashed] (h.south) -- ++(0,-0.8) node[below,
        font=\scriptsize\bfseries,text=alt]{живёт, пока жива лямбда};
    \end{tikzpicture}
  \end{center}

  \slidegap
  \aside{Захватили большой массив ради его длины и положили лямбду в
    долгоживущий объект --- массив \lo{не соберётся никогда}.
    \texttt{\hi{static}} перед лямбдой запрещает захват: компилятор
    проверит, и аллокации нет.}
\end{frame}

\begin{frame}[fragile]{События}
  \slidetop
\begin{lstlisting}[style=csharp,basicstyle=\ttfamily\small]
sealed class Downloader
{
    public event Action<int>? ProgressChanged;

    void Report(int percent) =>
        ProgressChanged?.Invoke(percent);
}

downloader.ProgressChanged += OnProgress;
downloader.ProgressChanged -= OnProgress;
\end{lstlisting}

  \slidegap
  \aside{Делегат, к которому снаружи можно только \hi{подписаться} и
    отписаться; вызывает владелец.}
\end{frame}

\begin{frame}[plain]
  \keyline{Делегат --- тип для метода, лямбда --- метод без имени,\\
    замыкание --- переменная, которая переехала в кучу.\\ Смотрите, что
    захватываете.}
\end{frame}

% ============================================================ 4. NULL ======
\section{4. Отсутствие значения}

\begin{frame}[plain]
  \sectionopener{4}{Отсутствие значения}{\texttt{null} --- самая дорогая
    ошибка в истории программирования.}
\end{frame}

\begin{frame}[fragile]{\texttt{int?}: значение и флаг}
  \slidetop
  \begin{minipage}[t]{0.5\textwidth}
\begin{lstlisting}[style=csharp,basicstyle=\ttfamily\small]
int? age = null;

age.HasValue              // False
age.GetValueOrDefault()   // 0
age + 1                   // null
age > 3                   // False
age.Value                 // exception
\end{lstlisting}
  \end{minipage}\hfill
  \begin{minipage}[t]{0.44\textwidth}
    \vspace{0.2em}
    \begin{cbul}{accent}
      \item \hi{отдельный тип}: структура из числа и флага <<есть ли
            значение>>
      \item арифметика с пустым даёт пустое
      \item сравнение с пустым --- ложь
    \end{cbul}
    \vspace{1.0em}
    \aside{Число не может быть \texttt{null}, а <<не указано>> бывает:
      возраст, дата. Для этого обёртка.}
  \end{minipage}
\end{frame}

\begin{frame}{Два механизма с одним знаком}
  \slidetop
  \versus{int? --- тип}{%
    \begin{cbul}{accent}
      \item \texttt{Nullable<int>}, структура
      \item существует во время выполнения
      \item \texttt{HasValue}, \texttt{Value}
    \end{cbul}}{string? --- аннотация}{%
    \begin{cbul}{alt}
      \item тот же \texttt{string}
      \item во время выполнения её нет
      \item читает только компилятор
    \end{cbul}}

  \slidegap
  \hrulethin{0.42}
  \rulegap
  \aside{Аннотация --- компромисс: двадцать лет кода считали, что ссылка может
    быть \texttt{null}. Новый тип сломал бы всё, аннотация не ломает ничего,
    а новый код получает \hi{предупреждения}.}
\end{frame}

\begin{frame}[fragile]{Компилятор думает о \texttt{null} вместе с вами}
  \slidetop
  \begin{minipage}[t]{0.42\textwidth}
\begin{lstlisting}[style=xml,basicstyle=\ttfamily\small]
<Nullable>enable</Nullable>
\end{lstlisting}
    \vspace{0.6em}
    \aside{Одна строка в \texttt{csproj} из Л1. В новых проектах уже
      включена.}
  \end{minipage}\hfill
  \begin{minipage}[t]{0.54\textwidth}
\begin{lstlisting}[style=csharp,basicstyle=\ttfamily\small]
static int Length(string? text)
{
    if (text is null)
    {
        return 0;
    }
    return text.Length;   // no warning
}
\end{lstlisting}
  \end{minipage}

  \slidegap
  \aside{Анализ потока понимает \texttt{return} после проверки,
    \texttt{\hi{is null}} и \texttt{\hi{is not null}}, оператор
    \texttt{??}. После проверки компилятор знает, что дальше не
    \texttt{null}, и молчит.}
\end{frame}

\begin{frame}[fragile]{Предупреждение ловит ошибку прода}
  \slidetop
\begin{lstlisting}[style=csharp,basicstyle=\ttfamily\small]
static string Greeting(User? user)
{
    var name = user?.Name ?? "guest";   // string? became string
    return $"Hello, {name}!";
}

static int NameLength(User user)
{
    return user.Name.Length;            // CS8602: Name may be null
}

sealed record User(string? Name);
\end{lstlisting}
  \vspace{0.4em}
  \aside{Второй метод компилируется, но с \lo{предупреждением} --- ровно там,
    где упал бы первый пользователь без имени.}
\end{frame}

\begin{frame}[fragile]{\texttt{?.} --- если есть}
  \slidetop
  \claim{Стратегия: \hi{распространить} отсутствие дальше.}

\begin{lstlisting}[style=csharp,basicstyle=\ttfamily\small]
var city = order?.Customer?.Address?.City ?? "unknown";
\end{lstlisting}
  \vspace{0.6em}
  \aside{Заказ, если есть, его покупатель, если есть, его адрес, если есть,
    город --- а если по дороге пусто, <<unknown>>. Старым способом --- четыре
    вложенных \texttt{if}.}

  \slidegap
  \hrulethin{0.42}
  \rulegap
  {\large \texttt{?.} \lo{прячет} отсутствие. Адрес обязан быть, а его нет
    --- строка молча вернёт <<unknown>>.}
\end{frame}

\begin{frame}[fragile]{\texttt{??} и \texttt{??=}}
  \slidetop
  \begin{minipage}[t]{0.52\textwidth}
    \colhead{alt}{через if}\par\vspace{0.5em}
\begin{lstlisting}[style=csharp,basicstyle=\twocol]
string name;
if (user.Name != null)
    name = user.Name;
else
    name = "guest";

if (_cache == null)
    _cache = new Dictionary<int, User>();
\end{lstlisting}
  \end{minipage}\hfill
  \begin{minipage}[t]{0.44\textwidth}
    \colhead{accent}{одной строкой}\par\vspace{0.5em}
\begin{lstlisting}[style=csharp,basicstyle=\twocol]
var name = user.Name ?? "guest";

_cache ??= new Dictionary<int, User>();
\end{lstlisting}
  \end{minipage}

  \slidegap
  {\renewcommand{\arraystretch}{1.4}
  \begin{tabular}{@{}p{1.6cm}p{11.6cm}@{}}
    \texttt{\hi{??}}  & <<или>>: \hi{заменить} отсутствие осмысленным ---
                        имя или <<гость>>, список или пустой список \\
    \texttt{\hi{??=}} & <<если пусто, присвоить>>: \hi{ленивая}
                        инициализация, кэш создаётся при первом обращении \\
  \end{tabular}}
\end{frame}

\begin{frame}{Что здесь значит \texttt{null}?}
  \slidetop
  \claim{Один \texttt{null} --- три разные вещи. Решите, какая, \hi{до}
    того, как писать \texttt{??}.}

  {\renewcommand{\arraystretch}{1.5}
  \begin{tabular}{@{}p{4.4cm}>{\raggedright\arraybackslash}p{8.8cm}@{}}
    \textbf{нормальный исход}   & <<не найден>> --- \hi{\texttt{User?}} в
                                  сигнатуре, вызывающий обязан обработать \\
    \textbf{ошибка контракта}   & пустой идентификатор --- \lo{исключение} на
                                  границе, внутри без вопросов \\
    \textbf{<<не загружено>>}   & ошибка проектирования: объект не создаётся
                                  наполовину, \texttt{required} \\
    \textbf{коллекция}          & \hi{пустая}, а не \texttt{null}. Никогда \\
  \end{tabular}}

  \slidegap
  \aside{\texttt{null} не должен молча пересекать границу слоя.}
\end{frame}

\begin{frame}[plain]
  \statement{Каждый \lo{\texttt{!}} --- это долг.\\[0.3em] {\color{muted}
    Рядом с ним --- комментарий, почему \texttt{null} здесь невозможен.}}
\end{frame}

% =========================================================== 5. ERRORS =====
\section{5. Ошибки}

\begin{frame}[plain]
  \sectionopener{5}{Ошибки}{Исключение --- это \hi{исключение}, а не поток
    управления.}
\end{frame}

\begin{frame}[fragile]{Неверный ввод --- не исключение}
  \slidetop
  \begin{minipage}[t]{0.47\textwidth}
    \colhead{alt}{исключение}\par\vspace{0.5em}
\begin{lstlisting}[style=csharp,basicstyle=\ttfamily\scriptsize]
try
{
    var age = int.Parse(input);
    Save(age);
}
catch (FormatException)
{
    Console.WriteLine("Not a number");
}
\end{lstlisting}
  \end{minipage}\hfill
  \begin{minipage}[t]{0.47\textwidth}
    \colhead{accent}{возвращаемое значение}\par\vspace{0.5em}
\begin{lstlisting}[style=csharp,basicstyle=\ttfamily\scriptsize]
if (int.TryParse(input, out var age))
{
    Save(age);
}
else
{
    Console.WriteLine("Not a number");
}
\end{lstlisting}
  \end{minipage}

  \slidegap
  \aside{Исключение раскручивает стек, собирает трассировку и проходит все
    \texttt{finally} --- на \lo{порядки} дороже возврата. Но главное ---
    \hi{скрытый переход}: управление оказывается в \texttt{catch} в другом
    файле.}
\end{frame}

\begin{frame}{Два вида плохого входа}
  \slidetop
  \versus{ошибка сценария}{%
    \begin{cbul}{accent}
      \item клиент всё сделал правильно
      \item{} <<средств недостаточно>>
      \item \hi{возвращаемое значение}
    \end{cbul}}{ошибка контракта}{%
    \begin{cbul}{alt}
      \item клиент нарушил договор
      \item{} <<сумма минус пять>>
      \item \lo{исключение}, один раз на границе
    \end{cbul}}

  \slidegap
  \hrulethin{0.42}
  \rulegap
  \aside{Можно описать в терминах домена, и клиент может осмысленно
    отреагировать --- это результат. Чего по договору быть не может --- это
    исключение контракта.}
\end{frame}

\begin{frame}[fragile]{Механика}
  \slidetop
  \begin{minipage}[t]{0.52\textwidth}
\begin{lstlisting}[style=csharp,basicstyle=\ttfamily\scriptsize]
try
{
    Process(file);
}
catch (FileNotFoundException)
{
    UseDefaults();
}
catch (IOException e)
{
    Log(e);
    throw;
}
finally
{
    file.Close();
}
\end{lstlisting}
  \end{minipage}\hfill
  \begin{minipage}[t]{0.44\textwidth}
    \vspace{0.2em}
    \begin{cbul}{accent}
      \item \texttt{catch} сверху вниз: \hi{частные выше общих}, иначе
            ошибка компиляции \texttt{CS0160}
      \item \texttt{\hi{throw;}} --- дальше с исходной трассировкой
      \item \texttt{\lo{throw e;}} --- трассировка \lo{перезаписана}
            (\texttt{CA2200})
      \item \texttt{finally} выполняется \hi{всегда}, даже после
            \texttt{return}
    \end{cbul}
  \end{minipage}
\end{frame}

\begin{frame}[fragile]{Бросок, фильтр, цепочка}
  \slidetop
\begin{lstlisting}[style=csharp,basicstyle=\ttfamily\fontsize{6.8}{7.8}\selectfont]
static int ParsePort(string? raw)
{
    try
    {
        return int.Parse(raw ?? throw new ConfigurationException("PORT is not set"));
    }
    catch (FormatException e) when (raw is { Length: 0 })
    {
        throw new ConfigurationException("PORT is empty", e);
    }
    catch (FormatException e)
    {
        throw new ConfigurationException($"PORT has a wrong format: '{raw}'", e);
    }
}

sealed class ConfigurationException(string message, Exception? inner = null)
    : Exception(message, inner);
\end{lstlisting}
\end{frame}

\begin{frame}{Где ловить всё}
  \slidetop
  \claim{Ловите только то, что можете \hi{обработать}.}

  \begin{center}
    \begin{tikzpicture}[font=\small,
      fr/.style={draw=hairline,line width=0.8pt,fill=codebg,
                 minimum width=5.6cm,minimum height=0.62cm,inner sep=0.2em},
      ar/.style={-{Straight Barb[length=1.6mm]},muted,line width=0.6pt}]
      \node[fr,draw=accent,line width=1.2pt,fill=paper] (b) at (0,0)
        {\texttt{\hi{catch (Exception)}} --- граница};
      \foreach \t/\y [count=\i] in {OrderController/-0.85,OrderService/-1.7,
                                    PaymentClient/-2.55} {
        \node[fr] (f\i) at (0,\y) {\texttt{\t}};
        \node[alt,font=\scriptsize,anchor=west] at (3.0,\y)
          {$\times$\ \texttt{catch (Exception)}};
      }
      \draw[ar] (f3.west) -- ++(-0.6,0) |- (b.west);
      \node[muted,font=\scriptsize,anchor=east] at (-3.45,-1.3)
        {летит вверх};
    \end{tikzpicture}
  \end{center}

  \slidegap
  \aside{Граница: точка входа утилиты, обработчик запроса (Л8), фоновая
    задача. Ниже <<поймать всё>> \lo{запрещено}: поймали баг, отмену,
    нехватку памяти и оставили систему в неизвестном состоянии.}
\end{frame}

\begin{frame}[fragile]{Как читать исключение}
  \slidetop
\begin{lstlisting}[style=term,basicstyle=\ttfamily\fontsize{6.9}{8.2}\selectfont]
$ PORT=abc dotnet run
Unhandled exception. ConfigurationException: PORT has a wrong format: 'abc'
 ---> System.FormatException: The input string 'abc' was not in a correct format.
   at System.Number.ThrowFormatException[TChar](ReadOnlySpan`1 value)
   at System.Int32.Parse(String s)
   at Program.<<Main>$>g__ParsePort|0_0(String raw) in /home/user/HelloWorld/Program.cs:line 8
   --- End of inner exception stack trace ---
   at Program.<<Main>$>g__ParsePort|0_0(String raw) in /home/user/HelloWorld/Program.cs:line 16
   at Program.<Main>$(String[] args) in /home/user/HelloWorld/Program.cs:line 1
\end{lstlisting}

  \slidegap
  \aside{Тип, сообщение, вложенное после \texttt{-{}-->}, трассировка. Читать
    с \hi{самого внутреннего}: оно ближе всего к причине. Файл и строка ---
    только если рядом лежит \texttt{pdb} (Л1).}
\end{frame}

\begin{frame}[plain]
  \keyline{Ловите только то, что можете обработать.\\ Остальное пусть
    долетит до границы: она одна, и она для этого существует.}
\end{frame}

\end{document}
